%%%%%%%%%%%%%%%%%%%%%%%%%%%%%%%%%
% Abstract                      %
%%%%%%%%%%%%%%%%%%%%%%%%%%%%%%%%%

\section*{Abstract}
\addcontentsline{toc}{section}{Abstract}

\vspace{1.4cm}
\large
\noindent

BAKO Motors Monitoring System is an embedded IoT solution developed as part of the ISS~496 senior project at MedTech Mediterranean Institute of Technology. The project addresses critical gaps in BAKO Motors' electric vehicle platform, including the absence of real-time sensor monitoring beyond the battery, no vehicle-wide data acquisition interface, and insufficient diagnostic documentation. The system is built around an ESP32 microcontroller interfaced with a CAN bus (via MCP2515/TJA1051T), a suite of 12 external sensors (current, voltage, temperature, handbrake, GNSS), and a custom Python-based sensor emulator for testing. A real-time web dashboard streams live vehicle telemetry via WebSocket, while a cloud pipeline transmits data over GPRS (SIM800L) to a SQLite database hosted on a VPS. A custom PCB was designed in KiCad to house all components. At the midway stage, the project has achieved 90\% completion in embedded development and 80\% in system design and dashboard visualisation, with future work focused on GNSS tracking, GSM/GPRS connectivity, and full in-vehicle PCB installation.\\

\textbf{Keywords:} Electric Vehicle, CAN Bus, ESP32, BMS, IoT, Real-Time Monitoring, Telemetry, PCB Design, BAKO Motor.
